\documentclass[../../../include/open-logic-section]{subfiles}

\begin{document}

\olfileid{sth}{z}{milestone}
\olsection{$\Z^-$: یک نقطهٔ عطف}

در بخش بعد به~\stagesinf{} بازخواهیم گشت. اما با اصل موضوع نامتناهی به
نقطهٔ عطف مهمی رسیده‌ایم. اکنون همهٔ اصول موضوع لازم برای نظریهٔ
$\Zminus$ را در اختیار داریم. به‌تفصیل:

\begin{defn}
نظریهٔ~$\Zminus$ این اصول موضوع را دارد: گسترش‌مندی، اجتماع، زوج‌ها،
مجموعه‌های توانی، نامتناهی و همهٔ نمونه‌های طرحوارهٔ جداسازی.
\end{defn}

این نام نشان‌دهندهٔ نظریهٔ مجموعه‌های \emph{زرملو} است (\emph{منهای}
چیزی که بعداً به آن خواهیم رسید). زرملو سزاوار این افتخار است، زیرا
اساساً این نظریه را در~\citeyear{Zermelo1908Untersuchungen} صورت‌بندی
کرد.\footnote{برای نکته‌های جالب تاریخی و فنی، بنگرید به
\citet[Appendix A]{Potter2004}.}

این نظریه به‌اندازه‌ای نیرومند است که مقدار عظیمی از ریاضیات را در آن
انجام دهیم. به‌ویژه، \emph{باید} به~\olref[sfr][][]{part} بازگردید و خود
را قانع کنید که هر کاری که در آنجا به‌شیوه‌ای ساده‌انگارانه انجام دادیم،
در~$\Zminus$ به صورتی صوری‌تر شدنی است. (پس از آنکه مدتی این کار را انجام
دادید، شاید بخواهید جلوتر بروید و~\olref[sth][z][arbintersections]{sec}
را بخوانید.) بنابراین، از این پس و بی‌هیچ توضیح بیشتری، خود را در حال کار
در~$\Zminus$ (یا نظریه‌ای دست‌کم به همان قوت) خواهیم دانست.

\end{document}
